\section{Discussion}
\label{sec:discussion}

\paragraph{Evaluation practice.}
Same-seed comparisons should check the three practices in
Section~\ref{sec:repair}: isolate probes, pin kernels, and verify each
materially distinct execution path at the resolution claimed.
Stream sharing alone does not imply divergence: the tested Trainer
evaluation is configuration-specifically stream-disjoint from training
(Table~\ref{tab:hf}). Existing isolation primitives help, but their use must
be checked. Deterministic execution and repeat/state audits incur runtime
and storage costs; some platforms lack deterministic implementations.

\paragraph{Interpretation and scope.}
Desynchronization, variance changes, and selection optimism are distinct.
Observed max-minus-mean uplift can contain both systematic placement
effects and optimism; the Gaussian formula applies only under its stated
model. Omitting selection removes this source of directional optimism,
not trajectory displacement. The no-op corollary is an application of the
pairing condition, not an accuracy claim about selection or averaging.

Our evidence is single-process and academic-scale, not distributed training,
ImageNet, or production-LLM validation. The nanoGPT experiment checks a
frozen mean prediction on held-out seeds in one setting, not universal
Gaussian scaling, per-run prediction, or untouched-test transfer.
Static exposure remains distinct from dynamic evidence, and historical
score-level equality does not establish state-level identity.
Appendix~\ref{sec:extended-context} records the detailed scope and practical
limitations.
